% TOP_PROBLEM: {"id":30007564,"problem_number":"LOCAL-30007564","title":"Learning to sustain cooperation","category_id":21,"category":{"id":21,"name":"economics","display_name":"Economics","description":"Open economic theory statements, published research agendas, and proposed scoped specifications; question provenance and historical priority are qualified per record.","slug":"economics","order_index":21,"created_at":"2026-10-08T00:00:00Z"},"status":"research_question","external_url":null,"legacy_ids":[],"published":false,"statement":"Estimate how learning about off-path punishments and forgiveness affects cooperation in the specified repeated prisoner’s dilemma under imperfect monitoring.","economics":{"original_id":53,"field":"Game theory and strategic interaction","kind":"empirical","formalization_basis":"adapted_research_question","source_status":"research_agenda","priority_status":"earliest_found","priority_note":"This is the earliest source in which the relevant agenda was verified during this audit; absolute priority has not been established.","earliest_verified_reference":{"authors":["Drew Fudenberg","David K. Levine"],"title":"Whither Game Theory? Towards a Theory of Learning in Games","year":2016,"url":"https://economics.mit.edu/sites/default/files/2022-09/Whither%20Game%20Theory%20Towards%20a%20Theory%20of%20Learning%20in%20Games.pdf","doi":"10.1257/jep.30.4.151","locator":"Active Learning: Learning about Off-Path Play, pp.162–163; Cooperation in Infinitely Repeated Games, pp.163–165","evidence_summary":"The discussion identifies counterfactual learning about rewards and punishments as an obstacle to explaining cooperation through learning."},"current_reference":{"authors":["Drew Fudenberg","David K. Levine"],"title":"Whither Game Theory? Towards a Theory of Learning in Games","year":2016,"url":"https://economics.mit.edu/sites/default/files/2022-09/Whither%20Game%20Theory%20Towards%20a%20Theory%20of%20Learning%20in%20Games.pdf","doi":"10.1257/jep.30.4.151","locator":"Active Learning: Learning about Off-Path Play, pp.162–163; Cooperation in Infinitely Repeated Games, pp.163–165","evidence_summary":"The discussion identifies counterfactual learning about rewards and punishments as an obstacle to explaining cooperation through learning."},"formalization_note":"The prisoner’s-dilemma primitives, randomized feedback intervention, and causal estimand make the source’s off-path-learning agenda concrete."},"statusQualification":"Published question or research agenda verified at the cited locator. The mathematical specification is adapted for this catalogue and includes additional assumptions; source publication does not establish first-ever priority or certify this exact model remains unsolved. The prisoner’s-dilemma primitives, randomized feedback intervention, and causal estimand make the source’s off-path-learning agenda concrete.","statusReviewedAt":"2026-10-08"}
% FIRST_POSTED: 2026-10-08
% ENTRY_KIND: statement_only
% !TeX program = lualatex
\documentclass[11pt]{article}
\usepackage{fontspec}
\usepackage[a4paper,margin=25mm]{geometry}
\usepackage{amsmath,amssymb,mathtools}
\usepackage{xurl}
\usepackage[hidelinks]{hyperref}
\newcommand{\sref}[2]{\hyperlink{source:#1:#2}{[S#2]}}
\setlength{\emergencystretch}{3em}
\begin{document}
% problemId:30007564
\section*{Learning to sustain cooperation}\label{30007564}
\subsection{Definitions and mathematical statement}
Two players repeatedly choose cooperation \(C\) or defection \(D\). Stage payoffs are \(2\) each at \((C,C)\), \(1\) each at \((D,D)\), and \(3\) to the defector and \(0\) to the cooperator otherwise. Continuation occurs independently with probability \(\beta\in(0,1)\). Players observe noisy private signals of past actions, with error probability \(\rho\), and receive randomized feedback treatment \(Z\) giving additional information about an opponent's contingent response to a deviation. Let \(Y_t\) indicate mutual cooperation at period \(t\), and define the normalized discounted cooperation outcome
\[
 C(Z)=(1-\beta)E\left[\sum_{t\geq1}\beta^{t-1}Y_t\mid do(Z)\right].
\]
Estimate \(C(1)-C(0)\) and determine how it varies with \(\beta\), \(\rho\), and players' elicited beliefs about punishment and forgiveness. Specify a structural learning model in which experimentation changes beliefs about these off-path responses, and test its counterfactual predictions on new monitoring regimes. Randomization occurs by pair to accommodate strategic spillovers; treatment-induced changes in preferences must be distinguished from information changes. Outcome censoring from random termination requires appropriate weighting. The target is causal learning mechanisms and their generalizability. Established folk-theorem possibilities for fully sophisticated players do not answer this empirical question about learning credible cooperation under imperfect feedback.

\par\textbf{Formulation and scope.} The prisoner’s-dilemma primitives, randomized feedback intervention, and causal estimand make the source’s off-path-learning agenda concrete.

\par\textbf{Open-question provenance.} An explicit question or research agenda was verified at the source locator. This does not by itself certify that every later version remains unresolved \sref{30007564}{1}.

\par\textbf{Historical priority.} This is the earliest source in which the relevant agenda was verified during this audit; absolute priority has not been established.

\subsection{Short English statement}
Estimate how learning about off-path punishments and forgiveness affects cooperation in the specified repeated prisoner’s dilemma under imperfect monitoring.

\subsection{Sources}
\begin{itemize}\raggedright
\item[\textnormal{[S1]}]\hypertarget{source:30007564:1}{} Drew Fudenberg, David K. Levine. 2016. Whither Game Theory? Towards a Theory of Learning in Games. Active Learning: Learning about Off-Path Play, pp.162–163; Cooperation in Infinitely Repeated Games, pp.163–165.
\url{https://economics.mit.edu/sites/default/files/2022-09/Whither%20Game%20Theory%20Towards%20a%20Theory%20of%20Learning%20in%20Games.pdf}.
DOI: 10.1257/jep.30.4.151.
{\small\textit{Earliest verified question/agenda reference; historical priority qualified below. The discussion identifies counterfactual learning about rewards and punishments as an obstacle to explaining cooperation through learning.}}
\end{itemize}
\end{document}
